% Source: https://artofproblemsolving.com/community/c3h1334880_tough_math_probplz_help
% Stable identity: aops-topic-1334880-q01
\textbf{(a)} The $n$th term is $n$ A's, then $n+1$ B's, then $n$ A's (check: $n=1$ gives ABBA). It has $n+1$ B's and $2n+(n+1)=3n+1$ letters.

\textbf{(b)} The ratio is $\boxed{\dfrac{n+1}{3n+1}}$.

\textbf{(c)} $\dfrac{n+1}{3n+1}=\dfrac{35}{100}=\dfrac{7}{20}$ gives $20n+20=21n+7$, so $n=\boxed{13}$ (the $13$th term has $14$ B's among $40$ letters).

\textbf{(d)} $\dfrac{n+1}{3n+1}=\dfrac{1}{3}$ would give $3n+3=3n+1$, which is impossible. In fact $\dfrac{n+1}{3n+1}-\dfrac{1}{3}=\dfrac{2}{3(3n+1)}>0$, so every term has slightly more than one third B's. The answer is $\boxed{\text{no}}$.
